\documentclass{article}
\usepackage{amsmath,amssymb}
\usepackage[margin=1in]{geometry}
\usepackage{enumitem}
\setlist[enumerate]{itemsep=2pt,topsep=2pt}

\begin{document}
	
	\section*{Databricks Certified Machine Learning Associate: 100 High-Probability Questions}
	
	% ============================================================
	\subsection*{Domain 1: Databricks Machine Learning (38\%)}
	% ============================================================
	
	\begin{enumerate}[label=\textbf{Q\arabic*.}]
		
		\item \textbf{MLflow Tracking}
		A data scientist needs parameters, metrics, and the model artifact logged automatically for a scikit-learn experiment. Which approach is best?
		\begin{enumerate}[label=\Alph*.]
			\item Manually write values to a CSV.
			\item Use \texttt{mlflow.autolog()} or explicit \texttt{mlflow.log\_param} and \texttt{log\_metric} calls within a run.
			\item Print values to the notebook output.
			\item Save a screenshot of the experiment.
		\end{enumerate}
		
		\item \textbf{Model Registry}
		A model must be promoted from development to production with a clear, auditable lifecycle. What is the recommended Databricks approach?
		\begin{enumerate}[label=\Alph*.]
			\item Email the pickle file to the production team.
			\item Register the model in Unity Catalog with versions and an alias such as \texttt{@champion}.
			\item Rename the artifact folder.
			\item Copy the model file to a shared DBFS path.
		\end{enumerate}
		
		\item \textbf{Feature Store}
		A team wants to avoid training-serving skew. Which Databricks capability is designed for this?
		\begin{enumerate}[label=\Alph*.]
			\item A shared Excel spreadsheet of feature definitions.
			\item Unity Catalog Feature Store with point-in-time joins.
			\item Recomputing features separately in the serving code.
			\item Using only raw columns at inference.
		\end{enumerate}
		
		\item \textbf{AutoML}
		A data scientist wants a quick baseline on a labeled tabular dataset. Which Databricks feature is most appropriate?
		\begin{enumerate}[label=\Alph*.]
			\item Databricks AutoML.
			\item Building a custom neural network from scratch.
			\item Deploying an endpoint before training.
			\item Creating an online feature table only.
		\end{enumerate}
		
		\item \textbf{ML Runtime}
		What is a primary advantage of using a Databricks Runtime for Machine Learning cluster?
		\begin{enumerate}[label=\Alph*.]
			\item It includes common ML libraries and integrations preinstalled.
			\item It eliminates all package conflicts automatically.
			\item It removes the need for cluster permissions.
			\item It allows skipping feature engineering entirely.
		\end{enumerate}
		
		\item \textbf{Unity Catalog Registry}
		What is a key benefit of registering models in Unity Catalog rather than the workspace registry?
		\begin{enumerate}[label=\Alph*.]
			\item Models become accessible only to the original author.
			\item Models are governed, discoverable, and promotable across workspaces at the account level.
			\item Models no longer need a version number.
			\item Models are automatically deployed to serving endpoints.
		\end{enumerate}
		
		\item \textbf{MLflow Search}
		A data scientist has many MLflow runs. How can they programmatically find the best run by a metric?
		\begin{enumerate}[label=\Alph*.]
			\item Use \texttt{mlflow.search\_runs()} and sort the resulting DataFrame.
			\item Manually click through each run in the UI.
			\item Export all runs to a CSV and use Excel.
			\item Delete runs until only one remains.
		\end{enumerate}
		
		\item \textbf{Model Aliases}
		What replaced the older model stages (Staging, Production, Archived) in Unity Catalog Model Registry?
		\begin{enumerate}[label=\Alph*.]
			\item Numeric version numbers only.
			\item Aliases such as \texttt{champion} and \texttt{challenger}.
			\item Tags only.
			\item File system paths.
		\end{enumerate}
		
		\item \textbf{Online vs Offline Features}
		What is the primary difference between online and offline feature tables?
		\begin{enumerate}[label=\Alph*.]
			\item Online tables are for training; offline tables are for serving.
			\item Online tables serve low-latency inference; offline tables are for batch training and analysis.
			\item Online tables cannot be created in Unity Catalog.
			\item Offline tables do not support point-in-time joins.
		\end{enumerate}
		
		\item \textbf{Feature Engineering Client}
		Which class is used to create and manage feature tables in Unity Catalog?
		\begin{enumerate}[label=\Alph*.]
			\item \texttt{FeatureEngineeringClient}.
			\item \texttt{SparkSession}.
			\item \texttt{MLflowClient}.
			\item \texttt{Pipeline}.
		\end{enumerate}
		
		\item \textbf{AutoML Output}
		What useful artifact does Databricks AutoML generate that allows a team to review and customize the winning approach?
		\begin{enumerate}[label=\Alph*.]
			\item A screenshot of the leaderboard only.
			\item Generated notebooks for data exploration and the best trial.
			\item A compiled binary with no source code.
			\item A cluster event log.
		\end{enumerate}
		
		\item \textbf{MLOps Best Practice}
		Which setup best reflects a sound MLOps strategy?
		\begin{enumerate}[label=\Alph*.]
			\item One shared workspace where notebooks are edited directly in production.
			\item Separate development, staging, and production environments with versioned code and tracked artifacts.
			\item Manual model file copies between users.
			\item Training new models without keeping prior run metadata.
		\end{enumerate}
		
		\item \textbf{Reproducibility}
		Which practice most directly improves reproducibility in a Databricks ML workflow?
		\begin{enumerate}[label=\Alph*.]
			\item Keeping training code in source control and logging run parameters and artifacts with MLflow.
			\item Relying on notebook revision history alone.
			\item Letting each data scientist rename columns locally.
			\item Saving only final predictions.
		\end{enumerate}
		
		\item \textbf{Model Promotion Strategy}
		In which scenario is promoting code preferred over promoting models?
		\begin{enumerate}[label=\Alph*.]
			\item When the model is a simple scikit-learn model.
			\item When the training pipeline is complex and the model artifact is small.
			\item When the model artifact is large and training is expensive.
			\item When the team uses only AutoML.
		\end{enumerate}
		
		\item \textbf{MLflow UI}
		Where in the MLflow UI can a user find the notebook that was run to create an experiment run?
		\begin{enumerate}[label=\Alph*.]
			\item The ``Source'' link for the run.
			\item The ``Artifacts'' tab only.
			\item The cluster event log.
			\item The Databricks home page.
		\end{enumerate}
		
		\item \textbf{Feature Store Creation}
		Which method creates a feature table in Unity Catalog?
		\begin{enumerate}[label=\Alph*.]
			\item \texttt{fe.create\_table()}.
			\item \texttt{spark.sql("CREATE TABLE")}.
			\item \texttt{mlflow.log\_artifact()}.
			\item \texttt{dbutils.fs.mkdirs()}.
		\end{enumerate}
		
		\item \textbf{Point-in-Time Lookup}
		Why is point-in-time lookup important when creating training sets from a feature store?
		\begin{enumerate}[label=\Alph*.]
			\item It ensures training does not use feature values that were unavailable at prediction time.
			\item It speeds up the query.
			\item It automatically imputes missing values.
			\item It converts categorical features to numeric.
		\end{enumerate}
		
		\item \textbf{Model Tags}
		How can a tag be set on a registered model in Unity Catalog?
		\begin{enumerate}[label=\Alph*.]
			\item \texttt{mlflow.set\_model\_tag()}.
			\item \texttt{spark.conf.set()}.
			\item \texttt{dbutils.notebook.exit()}.
			\item \texttt{display()}.
		\end{enumerate}
		
		\item \textbf{MLflow Run Identification}
		A data scientist wants to identify the best run using the MLflow Client API. Which method is used?
		\begin{enumerate}[label=\Alph*.]
			\item \texttt{client.search\_runs()}.
			\item \texttt{client.list\_experiments()}.
			\item \texttt{client.get\_run()}.
			\item \texttt{client.delete\_run()}.
		\end{enumerate}
		
		\item \textbf{Champion Promotion}
		How is a challenger model promoted to champion using aliases?
		\begin{enumerate}[label=\Alph*.]
			\item Delete the old model and upload the new one.
			\item Move the \texttt{@champion} alias to the new model version.
			\item Rename the model.
			\item Create a new experiment.
		\end{enumerate}
		
		% ============================================================
		\subsection*{Domain 2: Model Development (31\%)}
		% ============================================================
		
		\item \textbf{Regression Metrics}
		A regression model has RMSE = 25, MAE = 20, and $R^2 = 0.85$. Which statements are correct? (Select two.)
		\begin{enumerate}[label=\Alph*.]
			\item The model captures 85\% of the variance in the target.
			\item The RMSE being higher than MAE suggests larger errors exist.
			\item The model is underfitting.
			\item MAE is 25.
		\end{enumerate}
		
		\item \textbf{Imbalanced Classification}
		A binary classifier reports 98\% accuracy on data with 2\% positives. What should be done?
		\begin{enumerate}[label=\Alph*.]
			\item Ship the model immediately.
			\item Evaluate precision, recall, F1, and PR-AUC; consider class weights or resampling on training data only.
			\item Increase training epochs.
			\item Remove the positive class.
		\end{enumerate}
		
		\item \textbf{One-Hot Encoding}
		When is one-hot encoding least appropriate?
		\begin{enumerate}[label=\Alph*.]
			\item When the categorical feature has low cardinality.
			\item When the categorical feature has extremely high cardinality.
			\item When using linear models.
			\item When the feature is a string.
		\end{enumerate}
		
		\item \textbf{Imputation}
		For a feature with a highly skewed distribution, which imputation strategy is generally preferred?
		\begin{enumerate}[label=\Alph*.]
			\item Mean.
			\item Median.
			\item Mode.
			\item Zero.
		\end{enumerate}
		
		\item \textbf{Log Transformation}
		When is a log transformation most appropriate?
		\begin{enumerate}[label=\Alph*.]
			\item When the feature is normally distributed.
			\item When the feature is right-skewed and spans several orders of magnitude.
			\item When the feature is categorical.
			\item When the target is binary.
		\end{enumerate}
		
		\item \textbf{Estimator vs Transformer}
		In Spark ML, what is the key difference between an Estimator and a Transformer?
		\begin{enumerate}[label=\Alph*.]
			\item An Estimator has a \texttt{fit()} method; a Transformer has a \texttt{transform()} method.
			\item An Estimator only transforms data; a Transformer fits models.
			\item They are identical.
			\item A Transformer cannot be used in a Pipeline.
		\end{enumerate}
		
		\item \textbf{Spark ML Pipeline}
		Why should preprocessing and model training be wrapped in a Spark ML Pipeline?
		\begin{enumerate}[label=\Alph*.]
			\item To ensure the same preprocessing is applied at inference as during training.
			\item To reduce the number of columns.
			\item To avoid using MLflow.
			\item To eliminate the need for a test set.
		\end{enumerate}
		
		\item \textbf{Hyperopt Parallelization}
		Which Hyperopt component distributes trials across a Spark cluster?
		\begin{enumerate}[label=\Alph*.]
			\item \texttt{SparkTrials}.
			\item \texttt{Trials}.
			\item \texttt{fmin}.
			\item \texttt{hp.uniform}.
		\end{enumerate}
		
		\item \textbf{Hyperopt Search Algorithm}
		Which search algorithm does Hyperopt use by default via \texttt{tpe.suggest}?
		\begin{enumerate}[label=\Alph*.]
			\item Grid search.
			\item Random search.
			\item Tree-structured Parzen Estimator (Bayesian optimization).
			\item Exhaustive search.
		\end{enumerate}
		
		\item \textbf{Cross-Validation vs Train-Validation Split}
		What is the main trade-off between cross-validation and a single train-validation split?
		\begin{enumerate}[label=\Alph*.]
			\item Cross-validation is faster.
			\item Cross-validation gives a more stable estimate but is more computationally expensive.
			\item Train-validation split always gives higher accuracy.
			\item They are identical in cost.
		\end{enumerate}
		
		\item \textbf{Classification Metrics}
		Which metric is most appropriate for a highly imbalanced binary classification task?
		\begin{enumerate}[label=\Alph*.]
			\item Accuracy.
			\item Area under the Precision-Recall curve (PR-AUC).
			\item Mean Absolute Error.
			\item $R^2$.
		\end{enumerate}
		
		\item \textbf{Random Split}
		Which Spark operation randomly splits a DataFrame into training and test sets?
		\begin{enumerate}[label=\Alph*.]
			\item \texttt{DataFrame.randomSplit()}.
			\item \texttt{TrainValidationSplit}.
			\item \texttt{CrossValidator}.
			\item \texttt{DataFrame.where()}.
		\end{enumerate}
		
		\item \textbf{StringIndexer}
		Before one-hot encoding a string categorical column in Spark ML, which stage is typically required?
		\begin{enumerate}[label=\Alph*.]
			\item \texttt{StringIndexer}.
			\item \texttt{VectorAssembler}.
			\item \texttt{StandardScaler}.
			\item \texttt{Imputer}.
		\end{enumerate}
		
		\item \textbf{VectorAssembler}
		What is the purpose of \texttt{VectorAssembler} in Spark ML?
		\begin{enumerate}[label=\Alph*.]
			\item To combine multiple feature columns into a single vector column.
			\item To impute missing values.
			\item To encode categorical variables.
			\item To split the data.
		\end{enumerate}
		
		\item \textbf{Feature Scaling}
		Which scaler transforms features to have zero mean and unit variance?
		\begin{enumerate}[label=\Alph*.]
			\item \texttt{MinMaxScaler}.
			\item \texttt{StandardScaler}.
			\item \texttt{MaxAbsScaler}.
			\item \texttt{Normalizer}.
		\end{enumerate}
		
		\item \textbf{Data Leakage}
		A data scientist standardizes features before splitting into train and test. What is the risk?
		\begin{enumerate}[label=\Alph*.]
			\item The test set statistics leak into training, giving an overly optimistic evaluation.
			\item The model will underfit.
			\item The features will become categorical.
			\item Nothing; this is best practice.
		\end{enumerate}
		
		\item \textbf{Pandas UDF}
		What is a benefit of vectorized pandas UDFs over standard PySpark UDFs?
		\begin{enumerate}[label=\Alph*.]
			\item They process data in batches rather than row by row, improving performance.
			\item They work only on single-node DataFrames.
			\item They require no Python.
			\item They automatically handle missing values.
		\end{enumerate}
		
		\item \textbf{Target Encoding}
		Target encoding is most useful for which situation?
		\begin{enumerate}[label=\Alph*.]
			\item Low-cardinality categorical features.
			\item High-cardinality categorical features.
			\item Continuous features with outliers.
			\item Binary targets only.
		\end{enumerate}
		
		\item \textbf{Overfitting}
		A model performs very well on training data but poorly on test data. What is this called?
		\begin{enumerate}[label=\Alph*.]
			\item Underfitting.
			\item Overfitting.
			\item Data leakage.
			\item Regularization.
		\end{enumerate}
		
		\item \textbf{Bias-Variance}
		High variance in a model typically corresponds to which problem?
		\begin{enumerate}[label=\Alph*.]
			\item Overfitting.
			\item Underfitting.
			\item Data imbalance.
			\item Feature scaling.
		\end{enumerate}
		
		\item \textbf{Median vs Mean Imputation}
		When is median imputation preferred over mean imputation?
		\begin{enumerate}[label=\Alph*.]
			\item When the feature is symmetric.
			\item When the feature has outliers or is skewed.
			\item When the feature is categorical.
			\item When there are no missing values.
		\end{enumerate}
		
		\item \textbf{Class Weights}
		How can class imbalance be mitigated without resampling?
		\begin{enumerate}[label=\Alph*.]
			\item Use class weights in the model training.
			\item Remove all minority samples.
			\item Duplicate the majority class.
			\item Ignore the problem.
		\end{enumerate}
		
		\item \textbf{Evaluation Metric Selection}
		For a regression task where large errors are especially costly, which metric is most appropriate?
		\begin{enumerate}[label=\Alph*.]
			\item RMSE.
			\item MAE.
			\item Accuracy.
			\item F1.
		\end{enumerate}
		
		\item \textbf{Spark ML Model Persistence}
		How is a fitted Spark ML model saved?
		\begin{enumerate}[label=\Alph*.]
			\item \texttt{model.write().overwrite().save(path)}.
			\item \texttt{mlflow.log\_artifact()}.
			\item \texttt{spark.sql("SAVE")}.
			\item \texttt{dbutils.fs.put()}.
		\end{enumerate}
		
		% ============================================================
		\subsection*{Domain 3: ML Workflows / Data Processing (19\%)}
		% ============================================================
		
		\item \textbf{Summary Statistics}
		Which command computes summary statistics on a Spark DataFrame?
		\begin{enumerate}[label=\Alph*.]
			\item \texttt{df.summary()}.
			\item \texttt{df.show()}.
			\item \texttt{df.printSchema()}.
			\item \texttt{df.count()}.
		\end{enumerate}
		
		\item \textbf{Outlier Removal}
		How can outliers be removed based on standard deviation?
		\begin{enumerate}[label=\Alph*.]
			\item Filter values beyond a chosen number of standard deviations from the mean.
			\item Remove all values above the median.
			\item Remove the top 50\% of values.
			\item Use \texttt{dropDuplicates()}.
		\end{enumerate}
		
		\item \textbf{IQR Outlier Removal}
		What is the IQR method for outlier detection?
		\begin{enumerate}[label=\Alph*.]
			\item Remove values below Q1 or above Q3.
			\item Remove values below $Q1 - 1.5 \times IQR$ or above $Q3 + 1.5 \times IQR$.
			\item Remove values beyond the mean.
			\item Remove values with $z$-score $> 1$.
		\end{enumerate}
		
		\item \textbf{Pandas API on Spark}
		A data scientist has a pandas notebook and needs to scale to big data with minimal refactoring. What is the best approach?
		\begin{enumerate}[label=\Alph*.]
			\item Rewrite everything in Scala.
			\item Refactor to use the pandas API on Spark.
			\item Use only Spark SQL.
			\item Sample the data to fit in memory.
		\end{enumerate}
		
		\item \textbf{Visualization}
		Which method is appropriate for comparing two continuous features?
		\begin{enumerate}[label=\Alph*.]
			\item A scatter plot.
			\item A bar chart of counts.
			\item A pie chart.
			\item A box plot of a single feature.
		\end{enumerate}
		
		\item \textbf{Categorical Comparison}
		Which visualization is appropriate for comparing two categorical features?
		\begin{enumerate}[label=\Alph*.]
			\item A grouped bar chart or contingency table heatmap.
			\item A scatter plot.
			\item A line chart.
			\item A histogram of a continuous variable.
		\end{enumerate}
		
		\item \textbf{Missing Value Strategy}
		For a categorical feature, which imputation strategy is most appropriate?
		\begin{enumerate}[label=\Alph*.]
			\item Mean.
			\item Median.
			\item Mode.
			\item Linear interpolation.
		\end{enumerate}
		
		\item \textbf{Distribution Comparison}
		Which method compares two continuous distributions?
		\begin{enumerate}[label=\Alph*.]
			\item Overlaid histograms or density plots.
			\item A single pie chart.
			\item A bar chart of category counts.
			\item A confusion matrix.
		\end{enumerate}
		
		\item \textbf{Feature Store Training Set}
		Which method creates a training set from a feature table with point-in-time lookups?
		\begin{enumerate}[label=\Alph*.]
			\item \texttt{fe.create\_training\_set()}.
			\item \texttt{spark.read.table()}.
			\item \texttt{mlflow.log\_model()}.
			\item \texttt{Pipeline.fit()}.
		\end{enumerate}
		
		\item \textbf{Online Table Serving}
		What is the primary purpose of an online feature table?
		\begin{enumerate}[label=\Alph*.]
			\item Low-latency feature lookup for real-time inference.
			\item Batch training.
			\item Long-term storage.
			\item Data visualization.
		\end{enumerate}
		
		\item \textbf{Spark DataFrame Random Split}
		Which method splits a DataFrame into train and test sets with a seed?
		\begin{enumerate}[label=\Alph*.]
			\item \texttt{df.randomSplit([0.8, 0.2], seed=42)}.
			\item \texttt{df.split(0.8)}.
			\item \texttt{TrainValidationSplit}.
			\item \texttt{CrossValidator}.
		\end{enumerate}
		
		\item \textbf{CrossValidator}
		What does \texttt{CrossValidator} perform in Spark ML?
		\begin{enumerate}[label=\Alph*.]
			\item $k$-fold cross-validation for model selection and hyperparameter tuning.
			\item A single train-test split.
			\item Feature scaling.
			\item Model deployment.
		\end{enumerate}
		
		\item \textbf{TrainValidationSplit}
		What does \texttt{TrainValidationSplit} do?
		\begin{enumerate}[label=\Alph*.]
			\item Splits data once into training and validation sets for tuning.
			\item Performs $k$-fold cross-validation.
			\item Deploys a model.
			\item Encodes categorical variables.
		\end{enumerate}
		
		\item \textbf{Log Scale}
		When should a log transformation be applied to a feature?
		\begin{enumerate}[label=\Alph*.]
			\item When the feature is right-skewed or spans orders of magnitude.
			\item When the feature is binary.
			\item When the feature has no variance.
			\item When the feature is already normally distributed.
		\end{enumerate}
		
		\item \textbf{Outlier Impact}
		Why can outliers be problematic for mean imputation?
		\begin{enumerate}[label=\Alph*.]
			\item The mean is sensitive to outliers, so the imputed value may be biased.
			\item Outliers make the mean undefined.
			\item The median is always worse.
			\item Outliers improve the mean.
		\end{enumerate}
		
		\item \textbf{Data Summary}
		Which command provides descriptive statistics including count, mean, std, min, and max?
		\begin{enumerate}[label=\Alph*.]
			\item \texttt{df.describe()}.
			\item \texttt{df.show()}.
			\item \texttt{df.columns}.
			\item \texttt{df.dtypes}.
		\end{enumerate}
		
		\item \textbf{Categorical Encoding}
		What is the main drawback of one-hot encoding for very high-cardinality features?
		\begin{enumerate}[label=\Alph*.]
			\item It creates a very large number of sparse columns.
			\item It loses information.
			\item It cannot be used with Spark ML.
			\item It requires the target variable.
		\end{enumerate}
		
		\item \textbf{Mode Imputation}
		When is mode imputation most appropriate?
		\begin{enumerate}[label=\Alph*.]
			\item For continuous features with outliers.
			\item For categorical features with missing values.
			\item For binary targets.
			\item For normally distributed continuous features.
		\end{enumerate}
		
		\item \textbf{Pandas API Benefit}
		What is a key benefit of pandas API on Spark?
		\begin{enumerate}[label=\Alph*.]
			\item It allows pandas users to scale to big data with minimal code changes.
			\item It eliminates the need for Spark.
			\item It only works on single-node clusters.
			\item It replaces MLflow.
		\end{enumerate}
		
		% ============================================================
		\subsection*{Domain 4: Model Deployment (12\%)}
		% ============================================================
		
		\item \textbf{Serving Endpoint}
		What is the purpose of a Databricks serving endpoint?
		\begin{enumerate}[label=\Alph*.]
			\item To provide a REST API for real-time model inference.
			\item To store training data.
			\item To run SQL queries.
			\item To manage cluster permissions.
		\end{enumerate}
		
		\item \textbf{Scale to Zero}
		An ML engineer wants a serving endpoint active only during working hours to minimize cost. What should be enabled?
		\begin{enumerate}[label=\Alph*.]
			\item Scale to Zero.
			\item Autoscaling only.
			\item A larger cluster.
			\item Model versioning.
		\end{enumerate}
		
		\item \textbf{Batch vs Real-Time}
		What is the difference between batch inference and real-time inference?
		\begin{enumerate}[label=\Alph*.]
			\item Batch processes large volumes on a schedule; real-time serves individual requests with low latency.
			\item Batch is always faster.
			\item Real-time cannot use MLflow models.
			\item They are identical.
		\end{enumerate}
		
		\item \textbf{Model Deployment Registration}
		Before deploying a model to a serving endpoint, what step is typically required?
		\begin{enumerate}[label=\Alph*.]
			\item Register the model in Unity Catalog.
			\item Delete all prior runs.
			\item Convert the model to CSV.
			\item Stop the cluster.
		\end{enumerate}
		
		\item \textbf{Traffic Splitting}
		What is the purpose of traffic splitting between serving endpoints?
		\begin{enumerate}[label=\Alph*.]
			\item To route a percentage of traffic to different model versions for A/B testing.
			\item To increase training speed.
			\item To reduce storage costs.
			\item To impute missing values.
		\end{enumerate}
		
		\item \textbf{Streaming Inference}
		Which Databricks capability supports streaming inference?
		\begin{enumerate}[label=\Alph*.]
			\item Delta Live Tables with model serving.
			\item Only batch inference.
			\item Only real-time endpoints.
			\item Only AutoML.
		\end{enumerate}
		
		\item \textbf{Model Serving Scaling}
		What does autoscaling for serving endpoints do?
		\begin{enumerate}[label=\Alph*.]
			\item Adjusts the number of serving instances based on traffic.
			\item Automatically trains new models.
			\item Changes the model version.
			\item Deletes old experiments.
		\end{enumerate}
		
		\item \textbf{Deployment Consistency}
		Why must feature engineering be identical between training and serving?
		\begin{enumerate}[label=\Alph*.]
			\item To avoid training-serving skew and ensure consistent predictions.
			\item To reduce storage costs.
			\item To speed up cluster startup.
			\item To avoid using Unity Catalog.
		\end{enumerate}
		
		\item \textbf{Model Version Deployment}
		How can a specific model version be served from a Unity Catalog registry?
		\begin{enumerate}[label=\Alph*.]
			\item Reference the model by its version or alias in the serving endpoint configuration.
			\item Export the model to a CSV file.
			\item Rename the model.
			\item Use only the latest version automatically.
		\end{enumerate}
		
		\item \textbf{Deployment Monitoring}
		What is a key aspect of monitoring a deployed model?
		\begin{enumerate}[label=\Alph*.]
			\item Tracking prediction drift and data drift over time.
			\item Only monitoring cluster CPU usage.
			\item Ignoring performance after deployment.
			\item Only monitoring training accuracy.
		\end{enumerate}
		
		% ============================================================
		\subsection*{Additional High-Frequency Recall Items}
		% ============================================================
		
		\item \textbf{MLflow Autologging}
		What does \texttt{mlflow.autolog()} do?
		\begin{enumerate}[label=\Alph*.]
			\item Automatically logs parameters, metrics, and models for supported libraries.
			\item Automatically deploys a model.
			\item Automatically creates a cluster.
			\item Automatically deletes old runs.
		\end{enumerate}
		
		\item \textbf{Artifact Logging}
		How is a file artifact logged to an MLflow run?
		\begin{enumerate}[label=\Alph*.]
			\item \texttt{mlflow.log\_artifact(local\_path)}.
			\item \texttt{mlflow.log\_metric()}.
			\item \texttt{mlflow.log\_param()}.
			\item \texttt{mlflow.start\_run()}.
		\end{enumerate}
		
		\item \textbf{Parameter Logging}
		Which function logs a hyperparameter to MLflow?
		\begin{enumerate}[label=\Alph*.]
			\item \texttt{mlflow.log\_param()}.
			\item \texttt{mlflow.log\_metric()}.
			\item \texttt{mlflow.log\_artifact()}.
			\item \texttt{mlflow.register\_model()}.
		\end{enumerate}
		
		\item \textbf{Metric Logging}
		Which function logs a numeric metric to MLflow?
		\begin{enumerate}[label=\Alph*.]
			\item \texttt{mlflow.log\_metric()}.
			\item \texttt{mlflow.log\_param()}.
			\item \texttt{mlflow.log\_artifact()}.
			\item \texttt{mlflow.set\_tag()}.
		\end{enumerate}
		
		\item \textbf{Start Run}
		Which function starts a new MLflow run?
		\begin{enumerate}[label=\Alph*.]
			\item \texttt{mlflow.start\_run()}.
			\item \texttt{mlflow.end\_run()}.
			\item \texttt{mlflow.search\_runs()}.
			\item \texttt{mlflow.log\_model()}.
		\end{enumerate}
		
		\item \textbf{Register Model}
		Which function registers a model in Unity Catalog?
		\begin{enumerate}[label=\Alph*.]
			\item \texttt{mlflow.register\_model()}.
			\item \texttt{mlflow.log\_model()}.
			\item \texttt{mlflow.save\_model()}.
			\item \texttt{mlflow.export\_model()}.
		\end{enumerate}
		
		\item \textbf{Spark ML Evaluator}
		Which class evaluates a binary classifier in Spark ML?
		\begin{enumerate}[label=\Alph*.]
			\item \texttt{BinaryClassificationEvaluator}.
			\item \texttt{RegressionEvaluator}.
			\item \texttt{MulticlassClassificationEvaluator} only.
			\item \texttt{ClusteringEvaluator}.
		\end{enumerate}
		
		\item \textbf{Regression Evaluator}
		Which metric is available in \texttt{RegressionEvaluator}?
		\begin{enumerate}[label=\Alph*.]
			\item RMSE.
			\item Precision.
			\item Recall.
			\item F1.
		\end{enumerate}
		
		\item \textbf{Multiclass Evaluator}
		Which metric is available in \texttt{MulticlassClassificationEvaluator}?
		\begin{enumerate}[label=\Alph*.]
			\item F1.
			\item RMSE.
			\item MAE.
			\item $R^2$.
		\end{enumerate}
		
		\item \textbf{Feature Importance}
		How can feature importance be accessed for a tree-based Spark ML model?
		\begin{enumerate}[label=\Alph*.]
			\item \texttt{model.featureImportances}.
			\item \texttt{model.coefficients}.
			\item \texttt{model.intercept}.
			\item \texttt{model.summary}.
		\end{enumerate}
		
		\item \textbf{Linear Model Coefficients}
		How are coefficients accessed for a linear regression model in Spark ML?
		\begin{enumerate}[label=\Alph*.]
			\item \texttt{model.coefficients}.
			\item \texttt{model.featureImportances}.
			\item \texttt{model.intercept} only.
			\item \texttt{model.summary} only.
		\end{enumerate}
		
		\item \textbf{Intercept}
		How is the intercept of a linear model accessed?
		\begin{enumerate}[label=\Alph*.]
			\item \texttt{model.intercept}.
			\item \texttt{model.coefficients}.
			\item \texttt{model.featureImportances}.
			\item \texttt{model.summary}.
		\end{enumerate}
		
		\item \textbf{StringIndexer Handle Invalid}
		What happens by default when \texttt{StringIndexer} encounters an unseen label at transform time?
		\begin{enumerate}[label=\Alph*.]
			\item It throws an error unless \texttt{handleInvalid} is set.
			\item It silently ignores the row.
			\item It maps to index 0.
			\item It maps to the most frequent label.
		\end{enumerate}
		
		\item \textbf{OneHotEncoder Input}
		What type of column does \texttt{OneHotEncoder} expect as input in Spark ML?
		\begin{enumerate}[label=\Alph*.]
			\item A numeric index column, typically from \texttt{StringIndexer}.
			\item A raw string column.
			\item A continuous feature.
			\item A vector column.
		\end{enumerate}
		
		\item \textbf{VectorAssembler Input}
		What types of columns can \texttt{VectorAssembler} accept?
		\begin{enumerate}[label=\Alph*.]
			\item Numeric, boolean, and vector columns.
			\item Only string columns.
			\item Only categorical columns.
			\item Only binary columns.
		\end{enumerate}
		
		\item \textbf{Spark ML PipelineModel}
		What does fitting a \texttt{Pipeline} produce?
		\begin{enumerate}[label=\Alph*.]
			\item A \texttt{PipelineModel} that can be used for transformation and prediction.
			\item A raw DataFrame.
			\item A cluster.
			\item An MLflow experiment.
		\end{enumerate}
		
		\item \textbf{Model Persistence Format}
		What format does Spark ML use to save models?
		\begin{enumerate}[label=\Alph*.]
			\item A directory containing metadata and Parquet files.
			\item A single CSV file.
			\item A pickle file only.
			\item A JSON file only.
		\end{enumerate}
		
		\item \textbf{MLflow Model Flavor}
		What is an MLflow model flavor?
		\begin{enumerate}[label=\Alph*.]
			\item A convention for saving and loading models from different libraries (e.g., sklearn, spark).
			\item A type of cluster.
			\item A visualization tool.
			\item A data source.
		\end{enumerate}
		
		\item \textbf{Unity Catalog Three-Level Namespace}
		What is the three-level namespace for a model in Unity Catalog?
		\begin{enumerate}[label=\Alph*.]
			\item \texttt{catalog.schema.model}.
			\item \texttt{workspace.model.version}.
			\item \texttt{database.table.column}.
			\item \texttt{user.model.run}.
		\end{enumerate}
		
		\item \textbf{Model Signature}
		What is the purpose of a model signature when logging with MLflow?
		\begin{enumerate}[label=\Alph*.]
			\item It defines the expected input and output schema for the model.
			\item It encrypts the model.
			\item It compresses the model.
			\item It deletes the model after logging.
		\end{enumerate}
		
	\end{enumerate}
	
\end{document}